\documentclass[a4paper,11pt]{report}

% ========================= PACOTES =============================
\usepackage[portuguese]{babel}
\usepackage[utf8]{inputenc}
\usepackage[T1]{fontenc}

\usepackage{graphicx}
\usepackage{xurl}
\usepackage{hyperref}
\usepackage{float}

\usepackage{amsmath}
\usepackage{amssymb}
\usepackage{array}

\usepackage{parskip}
\usepackage{listings}
\usepackage{xcolor}
\usepackage{tikz}
\usetikzlibrary{arrows.meta, positioning, shapes.geometric}

\usepackage[a4paper,margin=2.5cm]{geometry}
\usepackage{booktabs}

% ========================= HIPERLIGAÇÕES =============================
\hypersetup{
    colorlinks=true,
    linkcolor=black,
    urlcolor=blue,
    citecolor=black
}

% ========================= COMANDOS ÚTEIS =============================
\newcommand{\code}[1]{\texttt{#1}}

\newcommand{\blankpage}{%
    \newpage
    \thispagestyle{empty}
    \mbox{}
    \newpage
}

\setlength{\parskip}{0.35em}
\setlength{\parindent}{0pt}

% ========================= LISTINGS =============================
\definecolor{lightgray}{RGB}{245,245,245}
\definecolor{darkgreen}{RGB}{0,128,0}

\lstdefinelanguage{csharp}{
    basicstyle=\ttfamily\small,
    keywordstyle=\color{blue}\bfseries,
    commentstyle=\color{gray}\itshape,
    stringstyle=\color{red},
    morecomment=[l]{//},
    morecomment=[s]{/*}{*/},
    morekeywords={class,struct,void,int,float,double,bool,var,null,const,
        override,enum,public,private,protected,using,namespace,static,
        new,this,return,if,else,true,false,string,Vector3,Quaternion,
        MonoBehaviour,Transform,GameObject,Camera,CharacterController,
        RequireComponent,DefaultExecutionOrder,Header,SerializeField,
        Time,Mathf,Physics,Debug,Animator,LayerMask,RaycastHit,
        List,Key,Gamepad},
    breaklines=true,
    showstringspaces=false,
    tabsize=4,
}

\lstset{
    language=csharp,
    backgroundcolor=\color{lightgray},
    frame=single,
    breaklines=true,
    aboveskip=0.5em,
    belowskip=0.5em,
    numbers=left,
    numberstyle=\tiny\color{gray},
    numbersep=5pt,
}

% ========================= TIKZ =============================
\tikzset{
    comp/.style={rectangle, draw, rounded corners, minimum width=3.6cm,
        minimum height=1cm, align=center, font=\small},
    arrow/.style={-{Latex}, thick},
}

\begin{document}

% ============================= CAPA ===============================
\begin{titlepage}
    \centering

    \includegraphics[width=0.33\textwidth]{logo_ubi.jpg}\\[0.7cm]

    {\large Universidade da Beira Interior}\\[0.2cm]
    {\large Departamento de Informática}\\[0.6cm]

    {\large Licenciatura em Engenharia Informática}\\[0.2cm]
    {\large Unidade Curricular: Projeto Final}\\[1.0cm]

    {\Huge\bfseries PaleoLusa}\\[0.2cm]
    {\large Protótipo XR sobre o Jurássico da Lourinhã}\\[1.0cm]

    {\large Relatório Técnico}\\[0.6cm]

    {\large Autor:}\\[0.25cm]
    Daniel Cardoso (Nº 52218)\\[0.8cm]

    {\large Orientador:}\\[0.25cm]
    {\large\bfseries Professor João Alfredo F. F. Dias, UBI}\\[0.8cm]

    {\large Covilhã, Março de 2026}
\end{titlepage}

% ========================= AGRADECIMENTOS ==========================
\chapter*{Agradecimentos}
\addcontentsline{toc}{chapter}{Agradecimentos}

Em primeiro lugar, agradeço ao meu orientador pelo acompanhamento, pela
disponibilidade e pelas sugestões que contribuíram para a qualidade deste
projeto.

Agradeço também à Universidade da Beira Interior e ao Departamento de
Informática por proporcionarem as condições necessárias ao desenvolvimento
deste trabalho.

Um obrigado especial aos utilizadores que participaram nos testes informais
do protótipo e cujo \textit{feedback} foi essencial para identificar melhorias.

Por fim, o meu sincero agradecimento à família e amigos pelo apoio constante
ao longo de todo o percurso académico.

\blankpage

% =========================== RESUMO ===============================
\chapter*{Resumo}
\addcontentsline{toc}{chapter}{Resumo}

Este relatório descreve o desenvolvimento de um protótipo XR em Unity centrado
no Jurássico da Lourinhã, com foco em edutainment. A aplicação permite ao
utilizador explorar um ambiente virtual inspirado num cenário jurássico, observar
dinossauros animados e aceder a conteúdos educativos através de interfaces
diegéticas integradas no mundo 3D.

O protótipo inclui três dinossauros principais: \textit{Stegosaurus},
\textit{Pachycephalosaurus} e \textit{Velociraptor}. Cada um possui animações,
elementos de áudio, colisores simplificados e painéis educativos próprios.
A interação com os dinossauros permite apresentar diferentes tipos de informação,
incluindo dados gerais da espécie, comparação de escala com um ser humano e uma
visualização anatómica/``raio-X''.

O projeto privilegia uma abordagem imersiva, evitando menus tradicionais fixos à
câmara e optando por painéis espaciais que acompanham a posição do utilizador,
mantendo a legibilidade em contexto XR. Para além da componente visual e
interativa, foram consideradas questões de desempenho, escala, estabilidade
técnica e preparação para utilização em dispositivos de Realidade Virtual.

\blankpage

\tableofcontents

\blankpage

% ====================== LISTA DE FIGURAS ==========================
\listoffigures
\addcontentsline{toc}{chapter}{Lista de Figuras}

\blankpage

% ====================== LISTA DE TABELAS ==========================
\listoftables
\addcontentsline{toc}{chapter}{Lista de Tabelas}

\blankpage

% ====================== LISTA DE ACRÓNIMOS ========================
\chapter*{Lista de Acrónimos}
\addcontentsline{toc}{chapter}{Lista de Acrónimos}

\begin{tabular}{@{}lp{10cm}@{}}
\textbf{2D}  & Bidimensional (\textit{Two-Dimensional}) \\[4pt]
\textbf{3D}  & Tridimensional (\textit{Three-Dimensional}) \\[4pt]
\textbf{AR}  & Realidade Aumentada (\textit{Augmented Reality}) \\[4pt]
\textbf{LOD} & Nível de Detalhe (\textit{Level of Detail}) \\[4pt]
\textbf{UI}  & Interface de Utilizador (\textit{User Interface}) \\[4pt]
\textbf{UX}  & Experiência do Utilizador (\textit{User Experience}) \\[4pt]
\textbf{VR}  & Realidade Virtual (\textit{Virtual Reality}) \\[4pt]
\textbf{XR}  & Realidade Estendida (\textit{Extended Reality}) \\[4pt]
\textbf{XRI} & XR Interaction Toolkit (Unity) \\[4pt]
\textbf{XRF} & XR Framework (Unity) \\[4pt]
\end{tabular}

\blankpage

% ============= DECLARAÇÃO DE USO DE IA GENERATIVA ================
\chapter*{Declaração de Uso de IA Generativa}
\addcontentsline{toc}{chapter}{Declaração de Uso de IA Generativa}

No âmbito do desenvolvimento deste projeto e da redação deste relatório,
foram utilizadas ferramentas de Inteligência Artificial Generativa, em
particular o assistente \textbf{Claude} (Anthropic). O recurso a estas
ferramentas teve os seguintes propósitos:

\begin{itemize}
    \setlength\itemsep{0pt}
    \item \textbf{Apoio ao desenvolvimento de código C\#:} assistência na
          implementação e depuração de scripts Unity, nomeadamente os sistemas
          de deambulação dos dinossauros, interação VR e painéis educativos;
    \item \textbf{Resolução de erros técnicos:} diagnóstico de erros de
          compilação e de comportamento em runtime, com explicação das causas
          e das correções aplicadas;
    \item \textbf{Apoio à redação:} revisão de estrutura e linguagem em
          algumas secções deste relatório, mantendo sempre a responsabilidade
          intelectual e de conteúdo do autor.
\end{itemize}

Todas as decisões de arquitetura, design e conteúdo educativo foram tomadas
pelo autor. O código gerado com assistência de IA foi revisto, adaptado e
integrado pelo autor, que assume plena responsabilidade pelo resultado final.
O uso destas ferramentas foi transparente e complementar ao trabalho próprio,
não substituindo o processo de aprendizagem e desenvolvimento individual.

\blankpage

% ======================= INTRODUÇÃO ===============================
\chapter{Introdução}

\section{Descrição da Proposta}

O projeto \textit{PaleoLusa} consiste no desenvolvimento de uma aplicação XR
em Unity centrada no Jurássico da Lourinhã. O objetivo é combinar imersão e
aprendizagem, permitindo ao utilizador explorar um ambiente jurássico e
interagir com dinossauros em escala aproximada, através de interfaces educativas
integradas no mundo virtual.

A experiência é orientada para edutainment, procurando apresentar conteúdos
científicos de forma acessível, visual e interativa. Em vez de recorrer apenas
a texto estático ou menus tradicionais, o projeto explora interfaces diegéticas
junto dos dinossauros, reforçando a sensação de presença e a ligação entre o
conteúdo educativo e o objeto observado.

\section{Enquadramento e Estado da Arte}

Tecnologias XR (VR/AR) são usadas em museus e património científico para tornar
conteúdos complexos mais fáceis de compreender. Em paleontologia, isto é
especialmente útil porque permite mostrar aquilo que uma exposição tradicional
nem sempre consegue: \textbf{escala real}, \textbf{anatomia},
\textbf{movimento/comportamento} e \textbf{contexto ambiental} de animais extintos.

Neste projeto, o objetivo é criar um protótipo XR em Unity sobre o
\textbf{Jurássico da Lourinhã}, combinando imersão e aprendizagem. A
experiência foca-se em ganhos educativos imediatos: observação dos dinossauros
em ambiente 3D, comparação de escala com humanos e visualização anatómica
através de painéis do tipo ``raio-X''.

A literatura e a prática em XR para paleontologia e património indicam que o
sucesso destas experiências depende sobretudo do design:

\begin{itemize}
    \setlength\itemsep{0pt}
    \item \textbf{Acesso sem fricção}: a adoção pode ser limitada quando o
          utilizador precisa de muitos passos para começar;
    \item \textbf{Valor claro}: o efeito novidade não é suficiente; a
          experiência tem de mostrar algo que não se consegue ver ou fazer
          facilmente de outra forma;
    \item \textbf{Educação leve e contextual}: informação curta e acionada
          por interação funciona melhor do que texto longo permanente no ecrã;
    \item \textbf{Qualidade visual e animação}: representações fracas reduzem
          impacto e credibilidade.
\end{itemize}

Com base nestes pontos, o protótipo privilegia interações curtas e claras,
interfaces espaciais e um escopo controlado, focando a qualidade da experiência
e a aprendizagem imediata.

\section{Objetivos do Projeto}

Os principais objetivos do projeto são:

\begin{itemize}
    \setlength\itemsep{0pt}
    \item Desenvolver um protótipo XR em Unity com ambiente explorável;
    \item Integrar múltiplos dinossauros animados e ajustados à escala do
          cenário;
    \item Implementar interações educativas através de UI diegética;
    \item Criar painéis informativos com informação geral, comparação de
          escala e visualização anatómica;
    \item Preparar a experiência para utilização em Realidade Virtual;
    \item Otimizar o protótipo para garantir estabilidade e fluidez;
    \item Documentar e justificar as principais decisões técnicas e de design.
\end{itemize}

\section{Contribuições}

As principais contribuições deste projeto são:

\begin{itemize}
    \setlength\itemsep{0pt}
    \item \textbf{Protótipo XR funcional} sobre o Jurássico da Lourinhã,
          dirigido a contextos educativos e museológicos, executável em
          headset Meta Quest;
    \item \textbf{Sistema de UI diegética} com três painéis por dinossauro
          (Informação, Comparação de Escala e Raio-X), integrado no ambiente
          3D sem recorrer a menus tradicionais fixos à câmara;
    \item \textbf{Sistema de comportamento autónomo} dos dinossauros com
          deambulação física (\textit{CharacterController}), deteção de
          obstáculos omnidirecional e colisão entre animais sem escalada;
    \item \textbf{Integração de controlos VR} compatível com o Meta Quest
          real através de \code{UnityEngine.XR.InputDevices}, com
          \textit{fallback} para o XR Interaction Simulator;
    \item \textbf{Documentação técnica} das decisões de implementação,
          problemas encontrados e soluções adoptadas, servindo de referência
          para projetos semelhantes em Unity XR.
\end{itemize}

\section{Organização do Documento}

Este relatório encontra-se organizado da seguinte forma:

\begin{itemize}
    \setlength\itemsep{0pt}
    \item \textbf{Capítulo 1:} Introdução, enquadramento e objetivos;
    \item \textbf{Capítulo 2:} Metodologia, planeamento e ferramentas
          utilizadas;
    \item \textbf{Capítulo 3:} Implementação do ambiente, dinossauros,
          interação e UI;
    \item \textbf{Capítulo 4:} Avaliação e testes;
    \item \textbf{Capítulo 5:} Resultados e discussão;
    \item \textbf{Capítulo 6:} Conclusões e trabalho futuro.
\end{itemize}

\section{Síntese do Capítulo}

Este capítulo apresentou o enquadramento do projeto \textit{PaleoLusa},
situando-o no contexto das tecnologias XR aplicadas a conteúdo paleontológico
e educativo. Foram definidos os objetivos principais, resumida a literatura
de referência e descrita a estrutura do documento. O capítulo seguinte
descreve a metodologia de desenvolvimento e as ferramentas adoptadas.

\blankpage

% ======================= METODOLOGIA ===============================
\chapter{Metodologia e Planeamento}

\section{Introdução}

Este capítulo descreve a abordagem metodológica seguida no desenvolvimento
do protótipo, as ferramentas e tecnologias adoptadas e o cronograma geral
do projeto. A escolha de uma metodologia incremental permitiu validar
funcionalidades de forma progressiva e adaptar o plano às necessidades
que foram surgindo ao longo do desenvolvimento.

\section{Abordagem de Desenvolvimento}

O desenvolvimento do projeto segue uma abordagem incremental, organizada por
fases. Numa primeira etapa, o foco recaiu sobre investigação, curadoria
científica, seleção de referências visuais e prototipagem inicial de escala.
Seguidamente, procedeu-se à construção do ambiente, integração dos modelos 3D
e implementação das animações dos dinossauros.

Numa fase posterior, foram introduzidas as componentes educativas e de
interação, nomeadamente interfaces diegéticas com múltiplos painéis. Por fim,
o trabalho passou a focar-se em polimento, correção de erros, otimização para
Realidade Virtual e documentação do processo.

Esta abordagem permitiu validar funcionalidades de forma progressiva, reduzindo
o risco de concentrar o desenvolvimento apenas na fase final.

\section{Ferramentas e Tecnologias}

As principais ferramentas e tecnologias utilizadas são:

\begin{itemize}
    \setlength\itemsep{0pt}
    \item \textbf{Unity 6}, como motor principal de desenvolvimento;
    \item \textbf{C\#}, para scripting e controlo dos sistemas interativos;
    \item \textbf{XR Interaction Toolkit 3.4.0}, para interação em contexto
          de Realidade Virtual;
    \item \textbf{XR Interaction Simulator}, para testar interações sem
          necessidade constante do headset;
    \item \textbf{Blender}, para eventuais ajustes de escala, importação ou
          preparação de modelos 3D;
    \item \textbf{Git e GitHub}, para controlo de versões;
    \item \textbf{Git LFS}, para gestão de ficheiros de maior dimensão.
\end{itemize}

\section{Cronograma (Alto Nível)}

O cronograma do projeto encontra-se estruturado da seguinte forma:

\begin{itemize}
    \item \textbf{Fevereiro --- Investigação e Prototipagem de Escala:}
    estudo sobre turismo científico e XR, seleção dos dinossauros e do
    ambiente a representar, configuração inicial do Unity e criação de uma
    base de escala.

    \item \textbf{Março --- Construção do Ambiente e Integração 3D:}
    importação de modelos 3D, configuração de materiais, iluminação,
    animações e comportamentos básicos dos dinossauros.

    \item \textbf{Abril --- Narrativa Educativa e UI:}
    criação de painéis informativos, comparação de escala, visualização
    anatómica e integração da UI diegética.

    \item \textbf{Maio --- Otimização, Testes e Escrita:}
    testes de usabilidade, otimização de desempenho, recolha de feedback,
    correção de erros e redação dos capítulos finais.

    \item \textbf{Junho --- Entrega e Defesa:}
    entrega do projeto e do relatório final, seguida de apresentação e
    defesa pública.
\end{itemize}

\section{Síntese do Capítulo}

Este capítulo descreveu a metodologia incremental adoptada no desenvolvimento
do \textit{PaleoLusa}, as principais ferramentas e tecnologias utilizadas e
o cronograma de trabalho distribuído por fase. A abordagem por fases garantiu
entregas parciais validáveis e facilitou a identificação precoce de problemas.
O capítulo seguinte detalha as decisões de implementação tomadas em cada
componente do sistema.

\blankpage

% ======================= IMPLEMENTAÇÃO ===============================
\chapter{Implementação}

\section{Introdução}

Este capítulo detalha as decisões de implementação do protótipo, cobrindo
a arquitetura geral do sistema, a construção do cenário, a integração dos
dinossauros e todos os sistemas desenvolvidos em C\#: animação, deambulação,
painéis educativos, locomoção, controlos VR e áudio. Cada secção documenta
as escolhas técnicas efectuadas, os problemas encontrados e as soluções
adoptadas durante o desenvolvimento.

\section{Arquitetura Geral do Sistema}

O projeto está organizado num conjunto de scripts C\# independentes que
comunicam entre si através de referências directas e de membros estáticos.
A Figura~\ref{fig:arch} ilustra as principais componentes e as suas relações.

\begin{figure}[H]
\centering
\begin{tikzpicture}[node distance=1.2cm and 1.6cm]

    % Coluna esquerda --- XR
    \node[comp] (xro) {XR Origin\\(XR Rig)};
    \node[comp, below=of xro] (grav) {XRCharacterCollisionMover\\(Gravidade)};
    \node[comp, below=of grav] (move) {ContinuousMoveProvider\\(Locomotion)};
    \node[comp, below=of move] (vrin) {VRInputHandler};

    % Coluna central --- Dinossauros
    \node[comp, right=of xro] (cam) {Main Camera\\(TrackedPoseDriver)};
    \node[comp, right=of grav] (sw) {DinoPanelSwitcher\\(por dinossauro)};
    \node[comp, right=of move] (lk) {DinoPanelLookAtPlayer\\(por painel)};
    \node[comp, right=of vrin] (anim) {DinoAnimationCycle\\(por dinossauro)};

    % Coluna direita --- Movimento e Áudio
    \node[comp, right=of cam] (wander) {StegoWander / PachyWander\\/ RaptorWander};
    \node[comp, right=of sw] (dinoaud) {DinoAudio\\(sons periódicos)};
    \node[comp, right=of lk] (stateaud) {StateAudioBehaviour\\(sons por estado)};
    \node[comp, right=of anim] (ambient) {AmbientAudio\\(AudioSource 2D)};

    % Setas
    \draw[arrow] (xro) -- (cam);
    \draw[arrow] (xro) -- (grav);
    \draw[arrow] (xro) -- (move);
    \draw[arrow] (xro) -- (vrin);
    \draw[arrow] (vrin) -- (sw);
    \draw[arrow] (sw)  -- (lk);
    \draw[arrow] (sw)  -- (anim);
    \draw[arrow] (cam) -- (lk);
    \draw[arrow] (anim) -- (wander);
    \draw[arrow] (anim) -- (stateaud);

\end{tikzpicture}
\caption{Arquitetura simplificada dos sistemas do protótipo.}
\label{fig:arch}
\end{figure}

O \code{XR Origin (XR Rig)} é o prefab central fornecido pelo XR Interaction
Toolkit. Sobre ele assentam todos os sistemas de locomoção e o nó de câmara.
Os scripts de lógica de jogo (\code{DinoPanelSwitcher}, \code{DinoPanelLookAtPlayer},
\code{DinoAnimationCycle}, scripts de deambulação e áudio) estão associados a
cada dinossauro individualmente, permitindo uma estrutura modular.

\section{Cenário e Ambiente}

O ambiente virtual foi desenvolvido em Unity com o objetivo de representar uma
interpretação visual de um ecossistema jurássico. O cenário inclui terreno
tridimensional gerado por múltiplos tiles de \textit{Terrain} do Unity,
vegetação, rochas, zonas de relevo e iluminação ambiental, procurando evitar
a sensação de ``caixa vazia'' e reforçar a narrativa ambiental.

O terreno é composto por tiles de 1000 $\times$ 1000 unidades. A camada
\textit{Ground} foi atribuída ao \textit{TerrainCollider} de cada tile, sendo
essa camada usada para deteção de solo pelos sistemas de física. A posição de
início do jogador fica próxima de $(-355, -31, 335)$ em coordenadas mundiais,
numa zona onde a superfície do terreno tem uma altitude negativa relativa à
origem da cena.

Durante o desenvolvimento foram realizados ajustes de escala entre o cenário,
o jogador e os dinossauros. Esta etapa revelou-se essencial, uma vez que em XR
a perceção de tamanho e distância tem impacto direto na sensação de presença.
Foram também efetuadas alterações ao cenário para melhorar a performance e
reduzir a complexidade visual.

\section{Integração dos Dinossauros}

Foram integrados três dinossauros principais:

\begin{itemize}
    \setlength\itemsep{0pt}
    \item \textbf{Stegosaurus} --- herbívoro de grande porte, com placas
          dorsais características;
    \item \textbf{Pachycephalosaurus} --- dinossauro de cúpula craniana
          espessa, de médio porte;
    \item \textbf{Velociraptor} --- carnívoro ágil de pequeno porte.
\end{itemize}

Cada dinossauro foi configurado com modelos 3D adquiridos em pacotes de assets
da \textit{Ferocious Industries} (disponíveis na Unity Asset Store), que incluem
malhas de alta qualidade, \textit{rigs} de animação, \textit{Animator Controllers}
pré-configurados e \textit{AudioClips} de sons comportamentais organizados por
categoria (rugidos, latidos, grunhidos, reações de dor).

A integração implicou ajustes significativos de escala, uma vez que os modelos
importados utilizavam unidades de escala distintas das do cenário. Esta calibração
foi feita com referência à escala humana (representada pelo \textit{XR Origin}),
garantindo que a perceção de tamanho em VR corresponde à realidade paleontológica
das espécies representadas. O \textit{Stegosaurus}, por exemplo, atinge
aproximadamente 4 metros de altura ao nível da bacia, o que se traduz numa escala
claramente superior à do utilizador e contribui para o impacto visual da experiência.

Os \textit{Animator Controllers} originais foram mantidos na sua estrutura base
e expandidos com estados adicionais criados especificamente para o projeto
(\code{WALKSTEGO}, \code{WALKPACHY}, \code{WALKRAPTOR}), que servem como pontos
de entrada para o sistema de deambulação (ver Secção~\ref{sec:wander}).

Para a componente física, foram utilizados colliders simplificados,
nomeadamente \textit{Capsule Collider} e \textit{Box Collider}, posicionados
manualmente para cobrir o volume aproximado de cada dinossauro. Esta escolha
foi feita por motivos de desempenho e estabilidade, evitando o uso de
\textit{Mesh Colliders} complexos que seriam proibitivos em contexto VR standalone.
Cada dinossauro tem também um componente \code{DinoFollowGround} que aplica um
\textit{raycast} descendente para manter o modelo fixo à superfície do terreno,
compensando irregularidades topográficas.

\section{Sistema de Animação dos Dinossauros}

Para evitar que os dinossauros permaneçam permanentemente no mesmo estado de
animação, foi implementado o script \code{DinoAnimationCycle}. Este script
alterna periodicamente entre animações configuradas no \textit{Animator} de
cada dinossauro, tornando o comportamento mais natural e dinâmico.

\begin{lstlisting}[caption={Excerto de DinoAnimationCycle.cs --- ciclo de animação.},
                   label={lst:animcycle}]
using UnityEngine;

public class DinoAnimationCycle : MonoBehaviour
{
    public Animator animator;

    [System.Serializable]
    public struct AnimState
    {
        public string stateName;
        [Min(0.1f)] public float duration;
    }

    public AnimState[] states;
    [Min(0f)] public float crossFadeTime = 0.5f;

    private int   _index;
    private float _timer;

    void Start()
    {
        if (animator == null)
            animator = GetComponentInChildren<Animator>();
        if (animator != null && states.Length > 0)
            PlayCurrent();
    }

    void Update()
    {
        if (animator == null || states.Length < 2) return;
        if (animator.speed == 0f) return; // pausado pelo DinoPanelSwitcher

        _timer -= Time.deltaTime;
        if (_timer <= 0f)
        {
            _index = (_index + 1) % states.Length;
            PlayCurrent();
        }
    }

    private void PlayCurrent()
    {
        animator.CrossFadeInFixedTime(states[_index].stateName, crossFadeTime);
        _timer = states[_index].duration;
    }
}
\end{lstlisting}

Cada estado do ciclo tem nome e duração configuráveis no \textit{Inspector}.
O ciclo é suspenso automaticamente quando \code{animator.speed == 0f}, o que
acontece durante a interação educativa através do \code{DinoPanelSwitcher}
(ver Secção~\ref{sec:panels}).

\section{Sistema de Deambulação dos Dinossauros}
\label{sec:wander}

Para tornar os dinossauros mais convincentes e naturais, foi implementado um
sistema de deambulação autónoma. Cada dinossauro percorre aleatoriamente a área
de jogo, alternando entre estados de repouso e caminhada de forma independente.

A solução final utiliza scripts individuais por espécie --- \code{StegoWander},
\code{PachyWander} e \code{RaptorWander} --- que comunicam com o \textit{Animator}
de cada dinossauro através de etiquetas de estado (\textit{tags}). Esta abordagem
foi necessária porque cada espécie tem uma estrutura de \textit{Animator} e de
\textit{root motion} distinta.

\subsection{Deteção do Estado de Caminhada por \textit{Tag}}

Em vez de controlar diretamente as transições do \textit{Animator}, os scripts
de deambulação observam passivamente o estado ativo e reagem a ele. Os estados de
caminhada em cada \textit{Animator} são marcados com a \textit{tag} \code{"Walk"}.
A cada \textit{frame}, o script verifica se o estado ativo tem essa \textit{tag}:

\begin{lstlisting}[caption={Deteção de estado de caminhada por tag (RaptorWander.cs).},
                   label={lst:wandertag}]
bool inWalkState = _animator.GetCurrentAnimatorStateInfo(0).IsTag("Walk");

if (inWalkState && !_walking)
    StartWalking();
else if (!inWalkState && _walking)
    StopWalking();
\end{lstlisting}

Esta abordagem desacopla o script de deambulação do \textit{Animator}: o
\code{DinoAnimationCycle} controla quando a caminhada começa e termina; o script
de deambulação apenas responde e aplica o deslocamento físico correspondente.
Quando a caminhada termina, o \code{DinoAnimationCycle} é reativado para retomar
o ciclo natural de animações.

\subsection{Movimento Físico com CharacterController}

O deslocamento é aplicado através do componente \code{CharacterController}
(atributo \code{[RequireComponent]} em cada script de deambulação), chamando
\code{\_cc.Move()} em vez de alterar diretamente \code{transform.position}.
Esta abordagem garante que os dinossauros são bloqueados por todos os colisores
da cena --- paredes, árvores e outros dinossauros --- sem atravessar objetos.
A gravidade é aplicada num método dedicado:

\begin{lstlisting}[caption={Gravidade aplicada via CharacterController (RaptorWander.cs).},
                   label={lst:wandergrav}]
void ApplyGravity()
{
    if (_cc.isGrounded)
        _verticalVelocity = -2f;
    else
        _verticalVelocity -= 9.81f * Time.deltaTime;

    _cc.Move(Vector3.up * _verticalVelocity * Time.deltaTime);
}
\end{lstlisting}

\subsection{Comportamento de Pausa e Desvio de Obstáculos}

Quando um obstáculo é detetado, o dinossauro não desvia imediatamente: para
por 1{,}5 a 3 segundos (\code{minWaitTime}/\code{maxWaitTime}) e escolhe de
seguida um novo alvo desviado entre 90° e 150° do eixo forward atual. Um
período de rotação (\code{\_steerCooldown} = 1{,}2 s) impede que o
\textit{SphereCast} sobreponha este giro com nova deteção prematura:

\begin{lstlisting}[caption={Pausa e escolha de direção lateral após obstáculo (RaptorWander.cs).},
                   label={lst:wanderwait}]
void StartWaiting()
{
    _waiting   = true;
    _waitTimer = Random.Range(minWaitTime, maxWaitTime);
}

void FinishWaiting()
{
    _waiting = false;
    Vector3 flatForward = new Vector3(
        transform.forward.x, 0f, transform.forward.z).normalized;
    float sideAngle = Random.value > 0.5f
        ? Random.Range(90f, 150f) : Random.Range(-150f, -90f);
    _target        = ClampToHome(transform.position +
        Quaternion.Euler(0f, sideAngle, 0f) * flatForward * wanderRadius);
    _steerCooldown = 1.2f;
}
\end{lstlisting}

A função \code{ClampToHome()} limita o alvo escolhido a um raio máximo
(\code{homeRadius} = 20 u) em torno da posição inicial de cada dinossauro,
evitando que os animais se afastem para fora da zona de jogo.

A deteção de obstáculos usa dois mecanismos complementares:

\begin{enumerate}
    \setlength\itemsep{0pt}
    \item \textbf{SphereCast frontal} --- cobre a direção de marcha e deteta
          árvores, paredes e outros obstáculos estáticos com antecedência;
    \item \textbf{OverlapSphere omnidirecional} --- verifica dinossauros
          em qualquer direção, resolvendo o facto de o evento
          \code{OnControllerColliderHit} não disparar de forma fiável
          em colisões entre dois \textit{CharacterControllers}:
\end{enumerate}

\begin{lstlisting}[caption={Deteção omnidirecional de outros dinossauros por OverlapSphere.},
                   label={lst:overlapsphere}]
int dinoMask = 1 << LayerMask.NameToLayer("Dino");
Collider[] nearbyDinos = Physics.OverlapSphere(
    transform.position + Vector3.up, 3f, dinoMask,
    QueryTriggerInteraction.Ignore);
foreach (var col in nearbyDinos)
{
    if (col.transform == transform ||
        col.transform.IsChildOf(transform)) continue;
    StartWaiting();
    return;
}
\end{lstlisting}

\subsection{Camada de Física \textit{Dino} e Prevenção de Escalada}

Todos os dinossauros foram atribuídos a uma camada de física dedicada
(\textit{Dino}). A \textit{Layer Collision Matrix} em
\textit{Project Settings $\to$ Physics} mantém \textit{Dino vs Dino}
ativado (colisão física entre dinossauros) mas os scripts de deambulação
definem \code{\_cc.stepOffset = 0.05f} em \code{Awake()}, reduzindo o
degrau automático do \textit{CharacterController} a 5 mm. Esta combinação
impede que um dinossauro escale o colisor em cápsula de outro ao colidirem
lateralmente, mantendo ao mesmo tempo a colisão física que evita que os
animais se atravessem mutuamente.

Os \textit{Near-Far Interactors} dos controladores VR têm o \textit{Raycast
Mask} e o \textit{Physics Layer Mask} configurados para incluir a camada
\textit{Dino}, garantindo que a interação com os painéis diegéticos funciona
mesmo com os dinossauros nesta camada dedicada.

\section{Sistema de Painéis Educativos}
\label{sec:panels}

Uma das funcionalidades centrais do projeto é o sistema de UI diegética.
Em vez de utilizar menus tradicionais fixos à câmara, a informação é
apresentada através de painéis tridimensionais integrados no ambiente e
associados diretamente a cada dinossauro.

Ao interagir com um dinossauro, surge um painel junto ao animal. Este painel
contém três modos principais, todos implementados na UI diegética:

\begin{itemize}
    \setlength\itemsep{0pt}
    \item \textbf{Informação}: dados gerais sobre a espécie, dieta, período
          e características principais;
    \item \textbf{Comparação}: relação de escala entre o dinossauro e um
          ser humano, com representação visual integrada no painel;
    \item \textbf{Raio-X}: representação anatómica do dinossauro através de
          imagem integrada no painel.
\end{itemize}

Os três painéis estão implementados e funcionais. O conteúdo visual final
(imagens de referência científica para os modos Comparação e Raio-X)
encontra-se em fase de preparação e integração.

O script \code{DinoPanelSwitcher} gere todo o ciclo de vida do painel de
cada dinossauro. Utiliza uma lista estática de instâncias para garantir que
apenas um painel pode estar aberto em simultâneo.

\begin{lstlisting}[caption={Excerto de DinoPanelSwitcher.cs --- abertura e fecho de painéis.},
                   label={lst:switcher}]
public class DinoPanelSwitcher : MonoBehaviour
{
    private static DinoPanelSwitcher openInstance;

    public GameObject dinoInfoRoot;
    public GameObject infoPanel;
    public GameObject comparePanel;
    public GameObject xrayPanel;
    public Animator animator;

    private int currentPanel;
    private bool isOpen;

    public void TogglePanel()
    {
        if (isOpen) ForceClose();
        else        ForceOpen();
    }

    public static void NextOpenPanel()
    {
        if (openInstance != null)
            openInstance.NextPanel();
    }

    private void ForceOpen()
    {
        if (openInstance != null && openInstance != this)
            openInstance.ForceClose();

        isOpen = true;
        openInstance = this;
        dinoInfoRoot.SetActive(true);
        ShowPanel(currentPanel);
        if (animator != null) animator.speed = 0f;
    }

    private void ForceClose()
    {
        isOpen = false;
        currentPanel = 0;
        dinoInfoRoot.SetActive(false);
        openInstance = null;
        if (animator != null) animator.speed = 1f;
    }

    private void ShowPanel(int index)
    {
        currentPanel = Mathf.Clamp(index, 0, 2);
        infoPanel.SetActive(currentPanel == 0);
        comparePanel.SetActive(currentPanel == 1);
        xrayPanel.SetActive(currentPanel == 2);
    }
}
\end{lstlisting}

\subsection{Barreira de Proximidade}

Quando o painel diegético está aberto, o utilizador pode inadvertidamente
aproximar-se em demasia do dinossauro, fazendo o texto do painel sair do
campo de visão por clipping ou por distância excessiva. Para resolver este
problema, cada dinossauro tem um \code{SphereCollider} filho
(\textit{ProximityWall}, raio 2{,}5 u) que é ativado em conjunto com o
painel e desativado quando o painel fecha. Assim, enquanto a informação
estiver visível, o \textit{CharacterController} do jogador é fisicamente
impedido de se aproximar mais do que a distância de leitura confortável.
O colisor mantém-se na camada \textit{Default} para interagir com o
\textit{CharacterController} do jogador, que também está na camada
\textit{Default}.

\begin{lstlisting}[caption={Ativação da barreira de proximidade em DinoPanelSwitcher.cs.},
                   label={lst:proxwall}]
[Header("Barreira de Proximidade")]
public SphereCollider proximityWall;

private void PauseDino(bool pause)
{
    // ... pausa do animator e scripts ...
    if (proximityWall != null)
        proximityWall.enabled = pause;
}
\end{lstlisting}

\subsection{Indicador Visual de Interagibilidade (Reticle)}

Para comunicar ao utilizador que está a apontar para um dinossauro
interagível, o \code{VRInputHandler} altera o material de uma pequena esfera
(\textit{reticle}) colocada à frente da câmara. Quando o \textit{Raycast}
atinge um objeto com \code{DinoPanelSwitcher}, o reticle muda para um
material dourado; caso contrário mantém o material neutro.

\subsection{Orientação do Painel para o Utilizador}

Os painéis seguem a orientação do utilizador através do script
\code{DinoPanelLookAtPlayer}. Quando o painel é ativado, regista a direção
horizontal da câmara do utilizador em relação ao dinossauro e mantém essa
posição durante toda a interação.

\begin{lstlisting}[caption={Excerto de DinoPanelLookAtPlayer.cs --- seguimento do utilizador.},
                   label={lst:look}]
public class DinoPanelLookAtPlayer : MonoBehaviour
{
    public Camera playerCamera;
    public Transform dinosaur;
    public float orbitDistance = 1.2f;
    public float heightOffset  = 0.8f;
    public float rotationSpeed = 10f;

    private bool    positionSet;
    private Vector3 _dinoOffset;

    void OnEnable() { positionSet = false; }

    void LateUpdate()
    {
        if (playerCamera == null || dinosaur == null) return;

        // Fixa a posicao uma vez ao abrir
        if (!positionSet)
        {
            Vector3 dir = Vector3.ProjectOnPlane(
                -playerCamera.transform.forward, Vector3.up);
            if (dir.sqrMagnitude > 0.001f)
            {
                _dinoOffset = dir.normalized * orbitDistance
                            + Vector3.up * heightOffset;
                positionSet = true;
            }
        }

        if (positionSet)
            transform.position = dinosaur.position + _dinoOffset;

        // Rotacao suave orientada para a camara
        Vector3 fullDir = -playerCamera.transform.forward;
        float hMag = new Vector2(fullDir.x, fullDir.z).magnitude;
        if (hMag > 0.001f)
        {
            float maxY = hMag * Mathf.Tan(40f * Mathf.Deg2Rad);
            fullDir.y = Mathf.Clamp(fullDir.y, -maxY, maxY);
            Quaternion target = Quaternion.LookRotation(
                fullDir.normalized) * Quaternion.Euler(0f, 180f, 0f);
            transform.rotation = Quaternion.Slerp(
                transform.rotation, target,
                Time.deltaTime * rotationSpeed);
        }
    }
}
\end{lstlisting}

\section{Sistema de Locomoção e Gravidade XR}
\label{sec:locomotion}

A locomoção do utilizador em XR é uma das componentes mais sensíveis do
projeto. O XR Interaction Toolkit 3.4.0 fornece o prefab
\textbf{XR Origin (XR Rig)}, que inclui:

\begin{itemize}
    \setlength\itemsep{0pt}
    \item Um \textit{CharacterController} com altura 1,8 m, centro em Y = 0,9 m
          e raio 0,25 m;
    \item Um \code{ContinuousMoveProvider} para locomoção contínua (WASD/joystick);
    \item Um \code{GravityProvider} para aplicar gravidade;
    \item Um \code{XRBodyTransformer} que integra todos os fornecedores de
          locomoção com o \textit{CharacterController}.
\end{itemize}

\subsection{GravityProvider}

O \code{GravityProvider} foi a solução escolhida para aplicar gravidade ao
jogador. Ao contrário de scripts personalizados que chamam diretamente
\code{CharacterController.Move()}, o \code{GravityProvider} processa o
movimento através do \code{XRBodyTransformer}, o que garante que a posição
da câmara simulada (via \code{TrackedPoseDriver}) é actualizada em
sincronismo com o corpo físico do jogador.

As definições mais relevantes do \code{GravityProvider} utilizadas no
projeto constam da Tabela~\ref{tab:gravity}.

\begin{table}[H]
\centering
\begin{tabular}{ll}
\toprule
\textbf{Parâmetro} & \textbf{Valor} \\
\midrule
Use Gravity                    & Ativado \\
Use Local Space Gravity        & Desativado \\
Update Character Controller Center & Desativado \\
Terminal Velocity              & 90 m/s \\
Gravity Acceleration Modifier  & 1 \\
Sphere Cast Layer Mask         & Ground \\
\bottomrule
\end{tabular}
\caption{Configuração do GravityProvider no XR Origin.}
\label{tab:gravity}
\end{table}

As definições \textbf{Use Local Space Gravity} e
\textbf{Update Character Controller Center} foram desativadas para evitar
conflitos com o \textit{XR Interaction Simulator} durante o desenvolvimento.
A camada \textit{Ground} no \textit{Sphere Cast Layer Mask} limita a deteção
de solo apenas ao \textit{TerrainCollider} do cenário, evitando falsos
positivos com outros colisores da cena.

\subsection{Delimitação da Área de Jogo}

Para restringir o jogador a uma zona específica do mapa, foram adotadas
duas abordagens complementares:

\textbf{Paredes invisíveis (BoundaryWalls):} foram criados objetos com
\textit{Box Collider} e \textit{MeshRenderer} desativado, posicionados
manualmente para delimitar o perímetro da área de jogo. Estas paredes
atuam sobre o \textit{CharacterController} do \textit{XR Origin} e bloqueiam
o avanço do jogador de forma impercetivelmente mas física.

\textbf{Vedação decorativa (PolygonPi):} sobre as paredes invisíveis foi
colocada uma vedação de madeira usando o asset \textit{PolygonPi Fence},
que suporta edição por \textit{spline} no editor Unity. O asset permite
desenhar livremente o traçado da vedação ao longo de qualquer curva,
combinando elemento visual com reforço da narrativa de que a zona de
observação se encontra delimitada.

O script \code{XRBoundsConstraint} permanece como segunda linha de defesa,
aplicando uma caixa de restrição AABB nos eixos X e Z para casos em que o
jogador consiga escapar às paredes físicas (por exemplo, durante teletransporte):

\begin{table}[H]
\centering
\begin{tabular}{ll}
\toprule
\textbf{Parâmetro} & \textbf{Valor} \\
\midrule
World Center & $(-355,\; -30{,}77,\; 335{,}52)$ \\
World Size   & $(800,\; 5,\; 800)$ \\
Constrain Y  & Desativado \\
Padding      & 0,25 m \\
\bottomrule
\end{tabular}
\caption{Configuração do XRBoundsConstraint (segunda linha de defesa).}
\label{tab:bounds}
\end{table}

\section{Mapeamento de Controlos VR}

O script \code{VRInputHandler} centraliza toda a lógica de entrada dos
controladores VR num único componente colocado no \textit{XR Origin}. Desta
forma evitam-se conflitos com os bindings internos do XR Interaction Toolkit
(locomoção, teleporte, agarrar).

A Tabela~\ref{tab:controls} resume o mapeamento implementado.

\begin{table}[H]
\centering
\begin{tabular}{lll}
\toprule
\textbf{Botão (controlador direito)} & \textbf{Ação XRI} & \textbf{Função} \\
\midrule
A & \textit{buttonSouth} & Abrir/fechar painel do dinossauro apontado \\
B & \textit{buttonEast}  & Ciclar sub-painel (Info $\to$ Comparação $\to$ Raio-X) \\
X & \textit{buttonWest}  & Fechar painel aberto \\
Y, gatilhos, grip & --- & Reservados para locomoção XRI \\
\bottomrule
\end{tabular}
\caption{Mapeamento de botões dos controladores VR.}
\label{tab:controls}
\end{table}

A interação com os dinossauros baseia-se num \textit{Raycast} a partir da
câmara. Quando o jogador prime o botão A a apontar para um dinossauro, o
script deteta o componente \code{DinoPanelSwitcher} no objeto atingido e
chama \code{TogglePanel()}. Se o raio não atingir nenhum dinossauro, o
painel aberto é fechado.

\begin{lstlisting}[caption={Excerto de VRInputHandler.cs --- leitura de botões XR e raycast.},
                   label={lst:vrinput}]
// Alias para resolver ambiguidade entre UnityEngine.InputSystem.InputDevice
// e UnityEngine.XR.InputDevice em tempo de compilacao
using XRInputDevice  = UnityEngine.XR.InputDevice;
using XRCommonUsages = UnityEngine.XR.CommonUsages;

// Obter dispositivos XR por no (funciona no headset real)
void RefreshDevices()
{
    if (!_rightHand.isValid)
    {
        var list = new List<XRInputDevice>();
        InputDevices.GetDevicesAtXRNode(XRNode.RightHand, list);
        if (list.Count > 0) _rightHand = list[0];
    }
}

bool GetXRButton(XRInputDevice device, InputFeatureUsage<bool> usage)
{
    if (!device.isValid) return false;
    device.TryGetFeatureValue(usage, out bool value);
    return value;
}

void Update()
{
    RefreshDevices();

    // Leitura dos botoes via XR InputDevices (Meta Quest real)
    bool aHeld = GetXRButton(_rightHand, XRCommonUsages.primaryButton);
    bool aPressed = !_prevA && aHeld;

    // Fallback para Gamepad (XR Interaction Simulator no editor)
    if (Gamepad.current != null)
        aPressed |= Gamepad.current.buttonSouth.wasPressedThisFrame;

    _prevA = aHeld;

    // Reticle: destaca quando aponta para dinossauro
    _currentPointed = GetPointedDino();
    if (reticle != null)
        reticle.material = _currentPointed != null
            ? reticleHighlight : reticleDefault;

    // Botao A --- abrir/fechar painel do dino apontado
    if (aPressed)
    {
        if (_currentPointed != null) _currentPointed.TogglePanel();
        else                         DinoPanelSwitcher.CloseOpenPanel();
    }
}

private DinoPanelSwitcher GetPointedDino()
{
    if (Physics.Raycast(_cam.transform.position,
                        _cam.transform.forward,
                        out RaycastHit hit, rayDistance))
        return hit.collider.GetComponentInParent<DinoPanelSwitcher>();
    return null;
}
\end{lstlisting}

A leitura de botões usa \code{UnityEngine.XR.InputDevices} como caminho
principal, com \textit{fallback} para \code{Gamepad} (XR Interaction Simulator
no editor). Foram necessários aliases de \code{using} para resolver a
ambiguidade de compilação entre \code{UnityEngine.InputSystem.InputDevice}
e \code{UnityEngine.XR.InputDevice}, que partilham o mesmo nome simples
quando ambos os \textit{namespaces} estão ativos.

\section{Desempenho e Otimização}

A otimização é um aspeto crítico em aplicações XR, uma vez que quebras de
desempenho podem causar desconforto ao utilizador. Durante o desenvolvimento
foram consideradas várias medidas:

\begin{itemize}
    \setlength\itemsep{0pt}
    \item Redução de objetos desnecessários no cenário;
    \item Utilização de colliders simplificados (\textit{Capsule}/\textit{Box})
          em vez de \textit{Mesh Collider};
    \item Utilização de materiais \textit{Unlit} em elementos de interface
          para evitar influência indesejada da iluminação;
    \item Preparação para \textit{baking} de iluminação;
    \item Separação entre elementos visuais e elementos interativos;
    \item Script de gravidade com instância única, verificada em cada frame
          para evitar cálculos redundantes.
\end{itemize}

Estas decisões procuram garantir maior estabilidade em dispositivos VR,
especialmente no Meta Quest, onde a performance é um fator determinante para
a qualidade da experiência.

\section{Sistema de Áudio}
\label{sec:audio}

O sistema de áudio foi estruturado em três camadas complementares, com o
objetivo de reforçar a imersão sem sobrecarregar a mistura sonora.

\subsection{Áudio Ambiente}

Um objeto vazio na cena contém um \code{AudioSource} configurado em modo 2D
(\textit{Spatial Blend} = 0), com \textit{Loop} e \textit{Play On Awake} ativos.
Este componente reproduz continuamente uma faixa de sons ambientes (vento, folhas,
água) a volume reduzido (0{,}3--0{,}4), garantindo que o utilizador nunca
experimenta silêncio total independentemente da sua posição no cenário.

\subsection{Sons Periódicos dos Dinossauros}

O script \code{DinoAudio} reproduz aleatoriamente sons de um conjunto configurável
de \textit{AudioClips}, com intervalos aleatórios entre valores mínimo e máximo
definidos no \textit{Inspector}:

\begin{lstlisting}[caption={Excerto de DinoAudio.cs --- sons periódicos.},
                   label={lst:dinoaudio}]
void Update()
{
    _timer -= Time.deltaTime;
    if (_timer <= 0f) { PlayRandom(); ScheduleNext(); }
}

void PlayRandom()
{
    if (clips == null || clips.Length == 0 || audioSource == null) return;
    if (audioSource.isPlaying) return;
    AudioClip clip = clips[Random.Range(0, clips.Length)];
    if (clip != null) audioSource.PlayOneShot(clip);
}
\end{lstlisting}

Cada dinossauro tem o seu próprio conjunto de clips (rugidos, grunhidos, latidos).
Sons de dor ou morte são excluídos desta camada por serem inadequados num contexto
de deambulação passiva; esses são reservados para a camada seguinte.

\subsection{Sons Sincronizados com Animações}

O script \code{StateAudioBehaviour}, derivado de \code{StateMachineBehaviour},
associa um \textit{AudioClip} diretamente a um estado do \textit{Animator}.
Quando o estado é ativado, o clip é reproduzido automaticamente:

\begin{lstlisting}[caption={StateAudioBehaviour.cs --- som acionado por estado do Animator.},
                   label={lst:stateaudio}]
public class StateAudioBehaviour : StateMachineBehaviour
{
    public AudioClip clip;
    [Range(0f, 1f)] public float volume = 1f;

    public override void OnStateEnter(Animator animator,
        AnimatorStateInfo stateInfo, int layerIndex)
    {
        if (clip == null) return;
        AudioSource src = animator.gameObject.GetComponent<AudioSource>();
        if (src != null) src.PlayOneShot(clip, volume);
    }
}
\end{lstlisting}

Este comportamento é aplicado a estados específicos como rugidos, ataques e
reações de dor, garantindo sincronização exata entre animação e som.

\subsection{Áudio 3D Espacial}

Os \code{AudioSource} dos dinossauros são configurados com \textit{Spatial Blend}
= 1 (totalmente 3D), \textit{Max Distance} = 15 m e \textit{Volume Rolloff} linear.
Este ajuste garante que os sons dos dinossauros dominam quando o utilizador está
próximo e diminuem gradualmente ao afastar-se, deixando o áudio ambiente em
evidência a longas distâncias. O \code{AudioListener} está posicionado na
\textit{Main Camera} (nó de câmara do \textit{XR Origin}), garantindo que o
cálculo de distância e direção do som é feito a partir da posição real da cabeça
do utilizador.

\section{Síntese do Capítulo}

Este capítulo documentou a implementação completa do protótipo \textit{PaleoLusa}.
As principais decisões técnicas foram a adopção de \code{CharacterController.Move()}
para deslocamento físico dos dinossauros (com pausa-e-desvio e \code{OverlapSphere}
omnidirecional), a camada de física \textit{Dino} para prevenir escalada entre
dinossauros, a UI diegética com painéis tridimensionais e barreira de proximidade,
o indicador visual de interagibilidade (\textit{reticle}), a delimitação de área
por paredes físicas e vedação decorativa, a reescrita do \code{VRInputHandler}
para compatibilidade com o Meta Quest real via \code{UnityEngine.XR.InputDevices},
e um sistema de áudio em três camadas complementares. O capítulo seguinte descreve
o protocolo de avaliação utilizado para validar a experiência com utilizadores reais.

\blankpage

% ======================= AVALIAÇÃO ===============================
\chapter{Avaliação}

\section{Introdução}

Este capítulo descreve o processo de avaliação do protótipo, incluindo o
protocolo de testes aplicado e as métricas utilizadas para medir a qualidade
da experiência. A avaliação centrou-se em sessões informais com utilizadores
sem experiência prévia em XR, com o objetivo de identificar problemas de
usabilidade e validar as decisões de design antes de uma versão mais completa.

\section{Protocolo de Testes}

A avaliação do protótipo foi realizada através de sessões de teste informais
com dois utilizadores sem experiência prévia em tecnologias XR. Cada sessão
decorreu individualmente, sem instrução prévia sobre os controlos ou objetivos,
de modo a simular a experiência de um utilizador novo. O ambiente foi executado
em modo editor com o \textit{XR Interaction Simulator}, que permite simular
a presença e os controlos de um headset VR sem necessidade do dispositivo físico.

A sessão foi estruturada em três fases:

\begin{enumerate}
    \setlength\itemsep{0pt}
    \item \textbf{Exploração livre} (aproximadamente 5 minutos): o utilizador
          percorreu o cenário sem orientação, com o objetivo de avaliar a
          navegabilidade espontânea e a primeira impressão do ambiente;
    \item \textbf{Tarefas guiadas}: o utilizador foi convidado a realizar
          um conjunto de tarefas específicas:
          \begin{itemize}
              \setlength\itemsep{0pt}
              \item Identificar e aproximar-se de pelo menos dois dinossauros;
              \item Abrir o painel educativo de um dinossauro;
              \item Navegar pelos três modos de painel (Informação, Comparação
                    e Raio-X);
              \item Fechar o painel e interagir com um segundo dinossauro;
          \end{itemize}
    \item \textbf{Entrevista breve pós-sessão}: recolha de impressões
          qualitativas orais sobre imersão, clareza da interface, interesse
          do conteúdo e sugestões de melhoria.
\end{enumerate}

As sessões não foram cronometradas de forma estrita, permitindo que cada
utilizador explorasse ao seu próprio ritmo. O observador registou notas sobre
dificuldades observadas, momentos de hesitação e reações espontâneas durante
a sessão.

\section{Métricas e Dimensões Avaliadas}

As dimensões avaliadas durante os testes incluem:

\begin{itemize}
    \setlength\itemsep{0pt}
    \item \textbf{Usabilidade}: facilidade de descoberta e uso dos controlos
          de interação (abrir/fechar/navegar painéis) sem instrução prévia;
    \item \textbf{Legibilidade}: clareza tipográfica e posicionamento espacial
          da informação nos painéis diegéticos;
    \item \textbf{Presença}: sensação subjetiva de estar num ambiente
          tridimensional habitado, influenciada pela qualidade visual, áudio
          e comportamento dos dinossauros;
    \item \textbf{Retenção de conhecimento}: capacidade de recordar pelo menos
          um facto sobre cada dinossauro observado após o final da sessão;
    \item \textbf{Interesse e envolvimento}: perceção subjetiva do utilizador
          sobre o interesse da experiência como ferramenta educativa;
    \item \textbf{Conforto}: ausência de desconforto, desorientação ou
          \textit{motion sickness} durante a navegação.
\end{itemize}

Dado o caráter informal e o número reduzido de participantes, os resultados
desta fase de avaliação têm natureza qualitativa e exploratória, servindo
principalmente para identificar problemas evidentes de usabilidade e validar
a direção geral do design antes da versão final.

\section{Síntese do Capítulo}

Este capítulo apresentou o protocolo de avaliação do protótipo, estruturado
em três fases (exploração livre, tarefas guiadas e entrevista pós-sessão),
e as dimensões de qualidade monitorizadas durante os testes. A avaliação
qualitativa foi concebida para identificar problemas evidentes e recolher
impressões genuínas de utilizadores novos. Os resultados detalhados são
apresentados no capítulo seguinte.

\blankpage

% ======================= RESULTADOS ===============================
\chapter{Resultados e Discussão}

\section{Introdução}

Este capítulo apresenta os resultados obtidos ao longo do desenvolvimento e
dos testes com utilizadores, discute as limitações identificadas no estado
atual do protótipo e analisa de que forma os resultados validam ou questionam
as decisões de design tomadas. O protótipo encontra-se numa fase funcional
avançada, com os três dinossauros animados, cenário explorável e sistema de
UI diegética com três painéis por dinossauro operacionais.

\section{Resultados Técnicos}

As principais funcionalidades implementadas e validadas são:

\begin{itemize}
    \setlength\itemsep{0pt}
    \item Integração funcional dos três dinossauros (\textit{Stegosaurus},
          \textit{Pachycephalosaurus}, \textit{Velociraptor}) no cenário;
    \item Ciclos de animação com alternância aleatória via \code{DinoAnimationCycle};
    \item Sistema de deambulação autónoma por espécie (\code{StegoWander},
          \code{PachyWander}, \code{RaptorWander}) baseado em
          \code{CharacterController.Move()}, com comportamento de pausa-e-desvio,
          \code{OverlapSphere} omnidirecional e território limitado por
          \code{homeRadius};
    \item Colisão física entre dinossauros sem escalada, via camada de física
          \textit{Dino} e \code{stepOffset = 0.05f} no \textit{CharacterController};
    \item Sistema de áudio completo: áudio ambiente 2D, sons periódicos por
          espécie (\code{DinoAudio}) e sons sincronizados com animações
          (\code{StateAudioBehaviour}) com áudio 3D espacial;
    \item Painéis educativos diegéticos implementados (Informação, Comparação
          e Raio-X) para cada dinossauro, com imagens de comparação de escala
          e raio-X integradas;
    \item Barreira de proximidade (\textit{ProximityWall}) que ativa ao abrir
          o painel e impede o jogador de se aproximar em demasia do dinossauro;
    \item Indicador visual (\textit{reticle}) que destaca em dourado quando o
          utilizador aponta para um dinossauro interagível;
    \item Sistema de interação VR reescrito para Meta Quest real via
          \code{UnityEngine.XR.InputDevices}, com \textit{fallback} para o
          XR Interaction Simulator;
    \item Gravidade e deteção de solo funcional via \code{GravityProvider};
    \item Delimitação de área de jogo com paredes físicas invisíveis, vedação
          decorativa (\textit{PolygonPi}) e \code{XRBoundsConstraint} como
          segunda linha de defesa;
    \item Eventos de animação dos assets de terceiros sem recetor resolvidos
          com guarda de limites de array em \code{StegaSoundEffects} e
          \code{BoneheadSoundEffects};
    \item Estrutura modular que permite adicionar novos dinossauros e painéis.
\end{itemize}

\section{Resultados dos Testes com Utilizadores}

Foram realizados dois testes informais com utilizadores adultos sem experiência
prévia em aplicações XR. Ambos os participantes conseguiram explorar o ambiente,
identificar os três dinossauros e abrir os painéis educativos sem qualquer
instrução prévia, o que valida a navegabilidade espontânea do protótipo.

Em termos de \textbf{usabilidade}, ambos os participantes descobriram a
interação com os dinossauros de forma autónoma, embora com alguma hesitação
inicial sobre qual o botão correto. Este resultado sugere a necessidade de um
indicador visual (por exemplo, um ícone ou destaque no dinossauro apontado)
que sinalize que o objeto é interagível antes de o utilizador pressionar o botão.

Relativamente à \textbf{legibilidade}, os textos dos painéis foram considerados
claros e bem dimensionados, com a informação suficientemente concisa para ser
lida sem esforço excessivo. O posicionamento lateral do painel em relação ao
dinossauro foi avaliado positivamente, uma vez que permite observar o animal
e consultar a informação simultaneamente.

Quanto à \textbf{presença e imersão}, ambos os participantes referiram
espontaneamente a qualidade visual do cenário e dos dinossauros como pontos
fortes da experiência. O comportamento autónomo dos dinossauros --- deambulação,
ciclos de animação e sons --- foi identificado como um fator relevante para a
sensação de estar num ambiente habitado.

Os principais \textbf{pontos de melhoria} identificados foram:
\begin{itemize}
    \setlength\itemsep{0pt}
    \item Ausência de indicações visuais claras sobre quais os objetos
          interagíveis e como iniciar a interação;
    \item Painéis de Comparação e Raio-X ainda sem conteúdo visual final,
          o que reduziu o impacto educativo dessas vistas;
    \item Desejo expresso de ouvir narração ou informação sonora sobre
          cada dinossauro, reforçando a pertinência do audioguia previsto
          como trabalho futuro.
\end{itemize}

\section{Limitações Identificadas}

As principais limitações no estado atual são:

\begin{itemize}
    \setlength\itemsep{0pt}
    \item Necessidade de otimização adicional para dispositivos VR standalone
          (Meta Quest): \textit{baking} de iluminação, LOD e redução de
          \textit{draw calls} ainda por concluir;
    \item Audioguia narrativo por espécie (narração informativa ao abrir o
          painel) não implementado --- o sistema de áudio atual cobre sons
          ambientes e comportamentais, mas não narração educativa;
    \item Testes em headset real (Meta Quest) não concluídos --- gravidade,
          áudio 3D, reticle e colisão do jogador devem ser validados
          diretamente no dispositivo;
    \item Conteúdo científico dos painéis sujeito a revisão de especialista
          em paleontologia antes de uma versão pública.
\end{itemize}

\section{Síntese do Capítulo}

Este capítulo apresentou os resultados técnicos obtidos e os resultados
dos testes informais com utilizadores, identificando os principais pontos
fortes do protótipo e as limitações que persistem. Os resultados confirmam
a adequação geral da abordagem de UI diegética e do sistema de comportamento
autónomo dos dinossauros, e apontam para as áreas prioritárias de
desenvolvimento futuro descritas no capítulo seguinte.

\blankpage

% ======================= CONCLUSÕES ===============================
\chapter{Conclusões e Trabalho Futuro}

\section{Introdução}

Este capítulo apresenta as conclusões do projeto \textit{PaleoLusa},
sintetizando os principais resultados alcançados e as aprendizagens retiradas
do processo de desenvolvimento. São também identificadas as perspetivas de
trabalho futuro mais relevantes para evoluir o protótipo para uma versão
completa e plenamente funcional.

\section{Conclusões}

O projeto \textit{PaleoLusa} apresenta um protótipo funcional com os
principais elementos necessários para uma experiência educativa XR:
cenário explorável, três dinossauros animados, interação VR e UI diegética
com múltiplos painéis educativos. Foram já realizados testes informais com
dois utilizadores, que permitiram validar a navegabilidade e a adequação
geral da experiência.

A escolha de painéis integrados no mundo virtual revelou-se adequada ao
contexto XR, uma vez que mantém a informação próxima do objeto observado e
evita quebrar a imersão com menus tradicionais. A integração de múltiplos
painéis por dinossauro permite organizar a informação de forma simples e
progressiva.

O desenvolvimento demonstrou a importância de decisões de design focadas na
experiência do utilizador, desempenho e estabilidade técnica, sobretudo em
aplicações destinadas a Realidade Virtual. Em particular, a integração correcta
dos sistemas de locomoção e gravidade do XR Interaction Toolkit exige atenção
especial ao fluxo interno do \code{XRBodyTransformer}, que deve ser o único
ponto de controlo do \textit{CharacterController} para evitar conflitos entre
sistemas.

\section{Trabalho Futuro}

Como trabalho a concluir e como perspetivas de desenvolvimento futuro, destacam-se:

\begin{itemize}
    \setlength\itemsep{0pt}
    \item Implementação do audioguia narrativo por espécie (narração
          informativa acionada ao abrir o painel educativo);
    \item Otimização para Meta Quest standalone: \textit{baking} de
          iluminação, LOD, redução de \textit{draw calls} e validação
          de framerate no dispositivo real;
    \item Testes em headset real (Meta Quest) e ajustes decorrentes;
    \item Revisão científica do conteúdo dos painéis por especialista
          em paleontologia;
    \item Expansão para novas espécies e novos pontos educativos;
    \item Avaliação formal com maior número de utilizadores.
\end{itemize}

\blankpage

% ======================= BIBLIOGRAFIA ===============================
\chapter*{Bibliografia}
\addcontentsline{toc}{chapter}{Bibliografia}

\begin{itemize}
    \item M. K. Bekele, R. Pierdicca, E. Frontoni, E. S. Malinverni, and J. Gain,
    ``A Survey of Augmented, Virtual, and Mixed Reality for Cultural Heritage,''
    \textit{Journal on Computing and Cultural Heritage}, vol. 11, no. 2,
    pp. 1--36, Mar. 2018.

    \item M. R. Bennett, M. Budka, M. Belvedere, and S. Reynolds,
    ``Deploying augmented reality to help engage the public with
    palaeontological applications,''
    \textit{Proceedings of the Geologists' Association}, vol. 133, no. 1,
    pp. 22--31, Feb. 2022.

    \item Unity Technologies. \textit{Unity Manual}. Disponível em:
    \url{https://docs.unity3d.com/Manual/}

    \item Unity Technologies. \textit{XR Interaction Toolkit Documentation}.
    Disponível em:
    \url{https://docs.unity3d.com/Packages/com.unity.xr.interaction.toolkit@latest}
\end{itemize}

\blankpage

% ======================= ASSETS UTILIZADOS =======================
\chapter*{Assets Utilizados}
\addcontentsline{toc}{chapter}{Assets Utilizados}

Esta secção lista os pacotes de assets de terceiros utilizados no desenvolvimento
do protótipo \textit{PaleoLusa}, adquiridos ou obtidos através da Unity Asset Store
ou de outras fontes.

\begin{itemize}

    \item Ferocious Industries. \textit{PBR Dinosaurs --- Velociraptor}.
    Unity Asset Store. Disponível em:
    \url{https://assetstore.unity.com}

    \item Ferocious Industries. \textit{PBR Dinosaurs --- Stegosaurus}.
    Unity Asset Store. Disponível em:
    \url{https://assetstore.unity.com}

    \item Ferocious Industries. \textit{PBR Dinosaurs --- Pachycephalosaurus}.
    Unity Asset Store. Disponível em:
    \url{https://assetstore.unity.com}

    \item PolygonPi. \textit{Fence Layout} (vedação por spline).
    Unity Asset Store. Disponível em:
    \url{https://assetstore.unity.com}

\end{itemize}

\end{document}
